%&pdflatex
\documentclass [12pt]{article}
\usepackage{authblk}
\usepackage{silence}
\usepackage{amsfonts}
\usepackage{amsmath}
\usepackage{wasysym}
\usepackage{hhline}
\usepackage{cite}
\usepackage{tocloft}
\usepackage[plainpages=false,pdfpagelabels]{hyperref}
\usepackage[nottoc]{tocbibind}

\WarningFilter{latexfont}{Font shape `U/wasy/b/n'}
\WarningFilter{latexfont}{Font shape `OMS/cmtt/m/n'}
\WarningFilter{latexfont}{Font shape `OMS/cmtt/bx/n'}

\setlength{\topmargin}{-2cm}
\setlength{\oddsidemargin}{0cm}
\setlength{\textheight}{23cm}
\setlength{\textwidth}{16cm}
\setlength{\parindent}{0pt}
\setlength{\parskip}{1.2pt}

\makeatletter
\renewcommand{\tt}{\begingroup\catcode`_=12\def\_{\string_}\tt@}
\newcommand{\tt@}[1]{\texttt{#1}\endgroup}
\newcommand{\ttsubsec}[1]{\subsection {\texorpdfstring{\tt{#1}}{#1}}}
\newcommand{\tr}{\mathop{\mathrm{tr}}}
\makeatother
\title
    {
    Gordion: Usage Notes
    \vspace*{-8pt}
    }
\date	{\small\today\vspace*{-30pt}}
\author	{Laurence G.~Yaffe}
\affil{%
    \small
    Department of Physics, University of Washington, Seattle WA 98195, USA
    \vspace*{-12pt}
    }%
\pagenumbering{roman}
\begin	{document}
\maketitle
\tableofcontents
%\newpage
\setlength{\parskip}{5pt plus 2pt minus 1pt}
\pagenumbering{arabic}

\section {Introduction}

\tt{Gordian} is a program for solving, numerically,
for low energy properties of
a variety U($N$) lattice gauge theories in the large $N$ limit using the
coherent state variational method.%
\footnote
    {%
    For background on this approach for solving large $N$ limits,
    see Refs.~\cite{YaffeRMP, CSVAI, CSVAII,Yaffe:2026ouz}.
    For information on the design of Gordion,
    see the [[design and implementation notes]].
    }
The name \tt{Gordian} is an allusion to the Greek legend
associated with Alexander the Great, who
is reputed to have sliced through a horrendously complicated
knot, instead of carefully untying it,
in Gordion (Latin: Gordium, Phrygian: Gordum),
the capital city of ancient Phrygia.
Dealing with arbitrarily complicated loops is at the
core of the coherent state variational method as
implemented in \tt{Gordian}.
It is hoped that this program helps cut through
the complexity of large $N$ gauge theories
{\large\smiley{}}.

Applicable theories are U($N$) lattice gauge theories,
in either Euclidean or Hamiltonian formulations,
on regular cubic lattices of dimension 1, 2, 3 or 4, with each
dimension either infinite or periodically compactified.
A fixed number of fundamental representation fermions
may also be present.%
\footnote
    {%
    A minimal staggered fermion lattice discretization
    is used with each fermion flavor having
    just $2N$ fermionic operators per lattice site,
    comprising one U($N$) fundamental representation
    plus its conjugate.
    }
All theories under consideration are invariant under cubic
lattice symmetries (translations, permutations and reflections)
together with charge conjugation (C) and time reversal,%
\footnote
    {%
    The Yang-Mills $\theta$ angle is zero.
    For Hamiltonian theories with fermions,
    only translations by an even number of lattice sites are
    required to be symmetries of the theory.
    }
and none of these symmetries are spontaneously broken.

\setlength{\parskip}{3pt}

Lattice Yang-Mills theories are defined by the following Hamiltonian
or Euclidean free energy:
\begin{subequations}
\begin{align}
	H_{\rm gauge}/N &= \tfrac 14 \lambda \sum_\ell \> \tr (E[\ell]^2)
	    + \lambda^{-1} \sum_p \> \tr (2 - U[\partial p] - U[\partial p]^\dagger) \,,
    \\
	F_{\rm gauge}/N &= \lambda^{-1} \sum_p \> \tr (2 - U[\partial p] - U[\partial p]^\dagger)
	    - S \,,
\end{align}
\end{subequations}
where $\ell$ labels positively oriented lattice links,
$p$ denotes positively oriented plaquettes,
$E[\ell]$ is the U($N$) electric field operator on link $\ell$, and
$U[\partial p]$ denotes the U($N$) gauge connection around plaquette $p$.
The relation between the dimensionless `t Hooft coupling $\lambda$ and a
conventionally defined continuum gauge coupling $g^2$ is
$\lambda \equiv g^2 N a^{4-d}$, with $a$ the lattice spacing and $d$ the spacetime dimension.

For QCD-like theories, the additional
fermion contributions to the Hamiltonian or free energy are:
\begin{subequations}\label{eq:fermidef}
\begin{align}
    H_{\rm fermi} &= \tfrac 12 \sum_{\langle xy \rangle} \>
	\psi^\dagger_x \, \eta_{\langle xy \rangle} \, U_{\langle xy \rangle} \, \psi_y
	+ m \, \lambda^p \sum_x \> \psi^\dagger_x \, \sigma_x \, \psi_x \,,
    \\
    F_{\rm fermi} &= \tfrac 12 \sum_{\langle xy \rangle} \>
	\psi^\dagger_x \, \eta_{\langle xy \rangle} \, U_{\langle xy \rangle} \, \psi_y
	+ m \, \lambda^p \sum_x \> \psi^\dagger_x \, \psi_x \,,
\end{align}
\end{subequations}
where $x$ labels lattice sites,
$\langle xy \rangle$ denotes the link connecting nearest-neighbor sites $x$ and $y$,
$U_{\langle xy \rangle}$ is the U($N$) gauge connection on link $\langle xy \rangle$,
and
$\eta_{\langle xy \rangle}$ is a non-dynamical $U(1)$ fermion flavor connection
satisfying $\eta[\partial p] = -1$ on every plaquette.
The lattice staggering factor
$\sigma_x$ equals $+1$ on the even sublattice and $-1$ on the odd sublattice.
For further details see Ref.~\cite{CSVAII}.

\section {Overall structure}

The key steps comprising the coherent state variation algorithm
\cite {CSVAII} are:
\begin{enumerate}\advance\itemsep -4pt
\item
    Construction of a truncated finite set of observables 
    which will be retained and whose expectation values will be computed.
    Here ``observables'' mean Wilson loops possibly containing
    (in Hamiltonian theories) one or two electric field insertions,
    or fermion bilinears possibly including one electric field insertion.
\item
    Selection of a finite set of coherence algebra generators
    which will be used to generate deformations of coherent states.
    These are anti-Hermitian combinations of Wilson loops,
    possibly with one electric field insertion, or fermion bilinears.
    Coefficients of these generators serve as Riemann normal coordinates
    in the neighborhood of any given point
    in the large $N$ phase space, and are used as variational parameters.
\item
    Evaluation, symbolically, of the commutators of all selected observables
    and generators.  These define geodesic equations
    on the large $N$ phase space and encode how observable expectation
    values vary as one deforms the state.
\item
    Evaluation, symbolically, of the commutators defining the first and second
    variations of the Hamiltonian (or Euclidean free energy) with respect
    to Riemann normal coordinates.
    Also, for Hamiltonian theories,
    evaluation of commutators of generators whose expectation values
    define the Lagrange bracket (the inverse Poisson bracket)
    of the large $N$ phase space.
\item
    Numerical minimization of the Hamiltonian (or Euclidean
    free energy), with each iterative step
    involving numerical evaluation of the symbolic expressions
    for the gradient and curvature of the Hamiltonian or free energy,
    solution of the linear equations predicting the location of the minimum,
    and numerical integration of the geodesic equations describing
    the change in expectation values of all observables as one moves
    to the predicted minimum.
\item
    In Hamiltonian theories,
    numerical evaluation of the symbolic expressions for the curvature of the
    Hamiltonian and the Lagrange bracket, in any chosen symmetry channel,
    and computation of the resulting
    small oscillation frequencies around the minimum,
    yielding a determination of the
    low-lying glueball or meson spectrum.
\end{enumerate}
%
Each of these steps is carried out first for the $O(N^2)$ pure gauge dynamics,
for some sequence of gauge coupling values,
and then again for the sub-leading $O(N)$ fundamental representation
fermion dynamics, for some sequence of fermion mass values.

Observable selection uses a classification based on the order in
a strong-coupling expansion at which an error is first introduced if
a given observable is omitted from an otherwise
correct strong-coupling expansion.%
\footnote
    {%
    For fermion bilinears, this is a combined strong-coupling
    and large mass expansion.
    For details, see Ref.~[[X]].
    }
For generators, the analogous classification is the order in a
strong-coupling expansion at which the generator is first needed.

\section {Typical usage}\label{sec:usage}

The program \tt{gordion} has a common code base covering all 
applicable theories, but differing executables are created for different theories.
So generating the correct executable proceeds actual calculations.
Using the standard \tt{Makefile},
with program source files residing in a \tt{Src} subdirectory, 
the command \tt{make THY=}\textit{theory}
will generate an executable named
\tt{gordion-}\textit{theory},
together with a \tt{gordion} symlink pointing to this executable.
Here \textit{theory} is one of the following designators
of a pre-defined target theory:%
\footnote
    {%
    Additional target theories with other lattice periodicities and/or fermion
    multiplicies may be easily defined by adding further cases in the source file
    \tt{Theory.h}.
    }
\begin{description}\advance\itemsep -3pt
\item [\tt {YM1e}]
    Euclidean Yang-Mills theory on a one-dimensional, period-one lattice,
    equivalent to the Euclidean one-plaquette model [[ref]];
\item [\tt {YM1h}]
    Hamiltonian Yang-Mills theory on a one-dimensional, period-one spatial lattice,
    equivalent to the Hamiltonian one-plaquette model [[ref]];
\item [\tt {YM2e}]
    Euclidean Yang-Mills theory on an infinite two-dimensional, lattice;
\item [\tt {YM2h}]
    Hamiltonian Yang-Mills theory on an infinite two-dimensional spatial lattice;
\item [\tt {YM3e}]
    Euclidean Yang-Mills theory on an infinite three-dimensional lattice;
\item [\tt {YM3h}]
    Hamiltonian Yang-Mills theory on an infinite three-dimensional spatial lattice;
\item [\tt {YM4e}]
    Euclidean Yang-Mills theory on an infinite four-dimensional lattice;
\item [\tt {QCD1e}]
    Euclidean QCD with one species of fermion on an infinite one-dimensional lattice
    (a free theory, included just for testing);
\item [\tt {QCD1h}]
    Hamiltonian QCD with one species of fermion on an infinite one-dimensional
    spatial lattice
    (a lattice formulation of the 1+1 dimensional `t Hooft model);
\item [\tt {QCD2e}]
    Euclidean QCD with one species of fermion on an infinite two-dimesional lattice;
\item [\tt {QCD2h}]
    Hamiltonian QCD with one species of fermion on an infinite two-dimesional 
    spatial lattice;
\item [\tt {QCD3e}]
    Euclidean QCD with one species of fermion on an infinite three-dimesional lattice;
\item [\tt {QCD3h}]
    Hamiltonian QCD with one species of fermion on an infinite three-dimesional 
    spatial lattice;
\item [\tt {QCD4e}]
    Euclidean QCD with one species of fermion on an infinite four-dimesional lattice;
\item [\tt {YM3ep}]
    Euclidean Yang-Mills theory on a three-dimensional lattice
    which is infinite in two dimensions and period one in the third dimension;
\item [\tt {YM3hp}]
    Hamiltonian Yang-Mills theory on a three-dimensional spatial lattice
    which is infinite in two dimensions and period one in the third dimension.
\end{description}

When the resulting executable \tt{gordion}-\textit{theory} is run, the program:
\begin{enumerate}
\item
    First loads any previously saved system information
    (\textit{sys-info} or \tt{.sys}) files or 
    vacuum expectation value (\textit{vev-data} or \tt{.vev}) files
    which are provided as command line arguments.
    These \textit{sys-info} or \textit{vev-data} files must
    be files previously created for the selected target theory
    (or for QCD theories, \textit{sys-info} or \textit{vev-data} files
    created for the corresponding Yang-Mills theory);
\item
    Then reads and executes commands from a command input file 
    specified with a ``\tt{-f }\textit{command-file}'' command line flag.
    If no such command file is specified, but a \tt{.gordion}
    file exists, this is assumed to be a default startup command file
    which is read and processed.
    (A ``\tt{-f -}'' command line flag suppresses reading of this
    default startup file.)
\item
    Subsequently reads and executes commands from standard input.
\end{enumerate}

The program implements a simple keyword-driven command line interpreter.
Commands may be terminated by either a semicolon (\tt{;}) or end of line.
The full set of reconized commands is listed in the following section.

As a concrete example, for a Yang-Mills theory
the following commands
illustrate the minimal steps of a calculation:
\begin{quote}
  \tt{build all 20}
\\\tt{do lambda 10 1 1}
\end{quote}
%
The \tt{build all 20} command tells the program to construct all observables
for the given target theory whose strong-coupling orders are less than or equal to 20.
For gauge observables on an infinite lattice, these are essentially all
Wilson loops constructed from up to five plaquettes, possibly
decorated (for Hamiltonian theories) by one or two electric field insertions.
Since no generators were explicitly defined,
following observable generation
the \tt{build all} command will install a minimal default
set of generators consisting, for Yang-Mills theory,
of one and two plaquette generators and, for QCD theories, also
one link fermion hopping terms.
Using these generators, the \tt{build} command then
performs the symbolic calculations needed to define the
gradient and curvature of the Hamiltonian (or free energy),
Lagrange bracket (for Hamiltonian theories),
and geodesic equations for all observables in the selected set.
Expressions for the curvature and Lagrange bracket are constructed
in all irreducible representations of the lattice cubic symmetry group (times C).

Given this information, the \tt{do lambda 10 1 1} command instructs
the program to minimize the Hamiltonian (or free energy) with the lattice
coupling $\lambda$ first set to 10, and then successively repeat
this minimization as $\lambda$ is decreased in steps of 1 down to 1.
After each minimization, the program writes a
human and Mathematica-readable results file.
For Hamiltonian theories, the results file contains
the minimized energy, small oscillation frequencies, and selected
observable expectation values for each value of $\lambda$. 
For Euclidean theories, the results file contains
the minimized free energy and selected observable expectation values.

Prior to the generation of symbolic expressions for the gradient, curvature,
Lagrange bracket and geodesic equations, one may define additional
generators (and their specified creation orders) with commands such as:
\begin{quote}
\tt{generator add (2) xyXY}
\\\tt{generator add (4) xxyXXY}
\\\tt{generator add (4) xyxYXyXY}
\\\tt{generator add (4) xyXYxyXY}
\\\tt{generator add (6) xyxYXyxYXyXY}
\\\tt{generator add (6) xyXYxyXYxyXY}
\\\tt{generator add (6) xxyXXYxyXY}
\\\tt{generator add (6) xxyyXYXY}
\\\tt{generator add (6) xxxyXXXY}
\\\tt{generator add (6) xxyXyxYXXY}
\\\tt{generator add (6) xxyxYXyXXY}
\\\tt{generator add (6) xyxyXYxYXyXY}
\\\tt{generator add (6) xyxYxyXYXyXY}
\end{quote}
Here, \tt{x} and \tt{y} denote link variables in the $+x$ and $+y$
lattice directions, while \tt{X} and \tt{Y} denote the corresponding
inverse links.
On a two dimensional lattice, the above commands define
the basic order 2 (one plaquette), 4 (two plaquette) and 6 (three plaquette)
Wilson loop generators.
With default settings, the program also installs corresponding generators 
containing one electric field insertion, produced by commuting the Wilson loop
generator with the gauge kinetic energy operator.
Every generator is automatically projected onto all possible
irreducible representations of the lattice cubic symmetry group (times C).

Command words, as well as subsequent argument words
other than explicit file, representation, or symmetry names,
may be highly abbreviated;
any initial matching sequence of one or more characters is sufficient
to match the first consistent target word.
For example, the above build command could be entered as just \tt{b all 20}.

The initial \tt{set autosave true} command instructs the program to
save the symbolic information produced during the build phase in
a \textit{sys-info} file, save the numerical observable expectation
values produced by the minimization phase (for each value of $\lambda$)
in a \textit{vev-data} file.
See below (under \tt{save sys} and \tt{save vev}) for details on the
naming and placement of these output files.

[[geoswap \& big calculations]]

[[fermions]]

\section {Command summary}\label{sec:commands}

The keyword-driven command line interpreter recognizes the following commands,
which may be terminated by either a semicolon (;) or end of line.
In the following summary,
\textit{italic} indicates user input,
brackets (\tt{[]}) indicate optional arguments,
and a vertical bar (\tt{|}) separates alternative arguments.
A plus mark$^+$ after the command indicates the command is only
applicable in Hamiltonian theories, while a double plus$^{++}$
indicates applicability only for Euclidean theories.

\begin{description}\advance\itemsep -1pt
\item [\tt \# {\it anything}]~\\
    A comment, to be ignored.

\item [\tt ?]~\\
    Print a command summary.

\item [\tt build all {\it max-order}]~\\
    Build observables (up to the specified maximum order)
    as well as the polynomial scripts for geodesics,
    gradient, curvature and (for Hamiltonian theories)
    the Lagrange bracket.

\item [\tt build curvature \tt{[}{\it repname}\tt{]}] \bgroup\itemsep 0pt \parskip 0pt
\item [\tt build curvature all]~\\
    Build the polynomial scripts
    used for evaluating the curvature of the
    Hamiltonian or free energy in the current representation.
    A specified representation becomes the current representation, while
    the final form builds the curvature scripts
    for all defined representations.
\egroup

\item [\tt build equations \tt{[}{\it repname}\tt{]}] \bgroup\itemsep 0pt \parskip 0pt
\item [\tt build equations all]~\\
    Build the polynomial scripts needed for evaluating the
    curvature and Lagrange bracket (for Hamiltonian theories)
    in the current repersentation and, if the representation is fully symmetric,
    also the Hamiltonian or free energy, gradient, and geodesic equations.
    A specified representation becomes the current representation, while
    the final form builds the curvature and Lagrange bracket scripts
    for all defined representations.
\egroup

\item [\tt build freeenergy$^{++}$]~\\
    Build the canonicalized observable polynomial for evaluating the
    free energy.

\item [\tt build geodesics]~\\
    Build the polynomial scripts
    used for evaluating the geodesic equations.

\item [\tt build gradient]~\\
    Build the polynomial scripts
    used for evaluating the gradient
    of the Hamiltonian or free energy. 

\item [\tt build hamiltonian$^+$]~\\
    Build the canonicalized observable polynomial for evaluating the
    hamiltonian.

\item [\tt build lagrange$^+$ \tt{[}{\it repname}\tt{]}] \bgroup\itemsep 0pt \parskip 0pt
\item [\tt build lagrange$^+$ all]~\\
    Build the polynomial scripts
    used for evaluating the 
    Lagrange bracket (for Hamiltonian theories) in
    the current representation.
    A specified representation becomes the current representation, while
    the final form builds the Lagrange bracket scripts
    for all representations.
\egroup

\item [\tt build observables {\it max-order}]~\\
    Build observables.
    All observables for the given minimization stage up to (and including)
    the specified maximum strong-coupling order are constructed.

\item [\tt call canon {\it obs}] \bgroup \itemsep 0pt \parskip 0pt
\item [\tt call classify {\it obs}]
\item [\tt call commute {\it op} {\it obs}]
\item [\tt call compose {\it symm$_1$} {\it symm$_2$}]
\item [\tt call do\_inner {\it Eloop-obs}]
\item [\tt call noEbound {\it obs}]
\item [\tt call transform {\it symm} {\it obs}]
\item [\tt call u1bound {\it obs}]~\\
    Call the indicated routines with the given arguments.
    This functionality is present for debugging purposes
    and is not needed for normal usage.
\egroup

\item [\tt do {\it coupling-name initial-value final-value increment}]~\\
    Evolve the specified coupling constant
    from the given initial value to the final value in steps of the
    given increment, performing Newton minimization at each step
    and writing results including the value of the Hamiltonian (or free energy),
    expectations of a selected subset of obseravbles and,
    for Hamiltonian theories, the small oscillation spectrum in all defined
    symmetry channels to an output \tt{.m} results file.

\item [\tt do minimize]~\\
    Perform a sequence of minimization steps,
    terminating when the $L_\infty$ norm of the predicted displacement from
    the minimum falls below the minimization tolerance (or when the
    limit on the maximum number of minimization steps is reached).

\item [\tt do step]~\\
    Perform one Newton minimization step,
    evaluating the gradient and curvature, solving for the
    predicted minimum, and integrating the geodesic equations
    to move to the predicted minimum.

\item [\tt evaluate curvature \tt{[}$*$\tt{]}]~\\
    Numerically evaluate the $T$-even curvature matrix
    produced by double commutators of $T$-even generators
    with the Hamiltonian (or free energy),
    using active generators in the currently selected irreducible representation.
    Without an optional \tt{*}, only the curvature eigenvalues are printed;
    with a \tt{*} the full $T$-even curvature matrix is first printed as well.

\item [\tt evaluate delta]~\\
    Evaluate, and print, the predicted location of the Hamiltonian or
    free energy minimum, equal to minus the inverse of the
    (vacuum representation) curvature times the gradient.

\item [\tt evaluate freeenergy$^{++}$]~\\
    Numerically evaluate the free energy and print the result.

\item [\tt evaluate geodesics {\it print-limit}] \bgroup \itemsep 0pt \parskip 0pt
\item [\tt evaluate geodesics $*$]~\\
    Compute the variation of all observables
    of the current minimization stage in the direction defined by \tt{delta}.
    Because this output vector can be very long, \tt{evaluate}
    \tt{geodesics} takes an argument specifying a limit on
    the number of observables whose derivatives will be printed,
    or else a \tt{*} requesting printing of the entire list.
\egroup

\item [\tt evaluate gradient]~\\
    Numerically evaluate, and print, the gradient of the Hamiltonian
    (or free energy) produced by all defined vacuum representation generators.

\item [\tt evaluate hamiltonian$^+$]~\\
    Numerically evaluate the Hamiltonian and print the result.

\item [\tt evaluate lagrange$^+$ \tt{[}$*$\tt{]}]~\\
    Evaluate the Lagrange bracket matrix,
    produced by expectations of generator commutators,
    for generators in the
    currently active irreducible representation.
    Without an optional \tt{*}, only the singular values of the Lagrange bracket
    matrix are printed;
    with a \tt{*} the full matrix is first printed as well.

\item [\tt evaluate metric$^+$ \tt{[}$*$\tt{]}]~\\
    Numerically evaluate the $T$-odd curvature matrix
    produced by double commutators of $T$-odd generator with the 
    Hamiltonian,
    using active generators in the
    currently selected irreducible representation.
    Without an option \tt{*}, the curvature eigenvalues are printed;
    with a \tt{*} the full $T$-odd curvature matrix is first printed as well.

\item [\tt evaluate polystats]~\\
    (Re)evaluates and prints polynomial script statistics,
    including the number of polynomial terms and their breakdown
    by monomial order for the gradient, curvature, Lagrange bracket
    (for Hamiltonian theories) and geodesic equations.
    (Only relevant for debugging or verifying sys-info file integrity.)

\item [\tt evaluate spectrum$^+$ \tt{[}{\it repname}\tt{]}] \bgroup\itemsep 0pt \parskip 0pt
\item [\tt evaluate spectrum$^+$ all]~\\
    Evaluate, and print, the
    small oscillation frequency (or mass) spectrum
    in the current symmetry representation
    derived from the eigenvalues of the generalized eigensystem
    of the curvature and Lagrange bracket.
    A specified representation becomes the current representation, while
    the final form evaluates the spectrum for all symmetry channels.
\egroup

\item [\tt generator add {\tt{[}}({\it order}){\tt{]}}
    {\tt{[}}{\it coeff$_1$}{\tt{]}} {\it op$_1$}
    {\tt{[[}}{\it coeff$_2$}{\tt{]}} {\it op$_2$}{\tt{]}} $\cdots$ ]~\\
    Construct the specified coherence group generator, 
    project it onto all defined irreducible representations,
    and add the results, when independent and non-vanishing, to the
    set of defined generators for each representation.
    One or more \tt{Op}'s (operators in the coherence algebra)
    may be given,
    each optionally preceded by a numerical coefficient;
    this is interpreted as a weighted sum of the indicated \tt{Op}'s.
    The generator order must either be specified explicitly
    as an initial parenthesized integer,
    or else the first given \tt{Op}, when canonicalized,
    must be a known observable.
    If multiple \tt{Op}'s are given and (when canonicalized)
    are known observables, they must all have a common creation order.

\item [\tt generator activate ({\it order})] \bgroup\itemsep 0pt \parskip 0pt
\item [\tt generator activate {\it{gen\#$_1$}} {\tt{[}}{\it gen\#$_2$}{\tt{]}} $\cdots$ ]~\\
    First form activates all generators of the given order, in all representations.
    Second form activates the numbered generator(s) of the currently active representation.
\egroup

\item [\tt generator suspend ({\it order})] \bgroup\itemsep 0pt \parskip 0pt
\item [\tt generator suspend {\it{gen\#$_1$}} {\tt{[}}{\it gen\#$_2$}{\tt{]}} $\cdots$ ]~\\
    First form suspends all generators of the given order, in all representations.
    Second form suspends the numbered generator(s) of the currently active representation.
\egroup

\item [\tt help]~\\
    Print a command summary.

\item [\tt load {\it sys-info-file}] \bgroup \itemsep 0pt \parskip 0pt
\item [\tt load \tt{[}{\it data-set\#}\tt{]} {\it vev-data-file}]~\\
    Load information from a previously saved compatible sys-info or vev-data file.
    If not specified, the default data set for vev-data files is the final one
    in the file.
    The theory type of the running program must match that of the specified file
    or, for QCD theories, match the parent Yang-Mills theory.
    The variable size (32 or 64 bit) used for observable indices and coefficients
    must also match.
    For vev-data files, the hash code(s) of the observable set(s)
    must match that of the corresponding sys-info file(s), and similarly
    the hash code for gauge observables in a QCD sys-info file must match
    that of the corresponding YM sys-info file.
    For sys-info files, if the file size is larger than half of system memory,
    the global \tt{geoswap} variable is set \tt{true} and the geodesic polynomial
    scripts are left on disk.
\egroup

\item [\tt load {\it data-set\#}] \bgroup \itemsep 0pt \parskip 0pt
\item [\tt load {\it coupling-name} {\it coupling-value}]~\\
    Load a specified data set from a currently open vev-data file.
    The first form uses the data-set number, as shown by the
    \tt{print vevindex} command.
    The second form compares the specified coupling value to the
    values of data records in the vev-file, and loads the last
    such matching record.
\egroup

\item [\tt print baseobs]~\\
    Print the list of basic observables (those used in the
    Hamiltonian or free energy).

\item [\tt print blablevels]~\\
    Print the defined verbosity levels.

\item [\tt print bucketlist]~\\
    Print the list of observable buckets used
    for multi-threading.

\item [\tt print canon]~\\
    Print the canonicalization cache.

\item [\tt print couplings]~\\
    Print the defined adjustable coupling constants.

\item [\tt print curvature $*$] \bgroup \itemsep 0pt \parskip 0pt
\item [\tt print curvature {\it term\#} {\tt{[}}{\it gen\#$_2$}
    {\tt{[}}{\it gen\#$_2$}{\tt{]}}{\tt{]}} ]~\\
    Print symbolic curvature expressions.
    The first form prints all curvature components.
    The second form allows one to print the curvature components
    of a specified term in the Hamiltonian or free energy
    (without any overall adjustable coupling constant)
    and, optionally, specify one or both
    indices of the generators in the active symmetry representation
    which act on the indicated term.
\egroup

\item [\tt print fermiinit]~\\
    Print the mapping from closed fermion bilinears to loops which
    is used for initialization of fermion bilinear expectation values.

\item [\tt print freeenergy$^{++}$]~\\
    Print the expression defining the free energy.

\item [\tt print generator] \bgroup \itemsep 0pt \parskip 0pt
\item [\tt print generator $*$]
\item [\tt print generator {\it gen\#}]~\\
    Print generator(s).
    The first form prints an abbreviated form of all generators in the
    current representation.
    The second form prints the full form of all generators in the
    current representation.
    The final form prints the full form of a particular generator.
    For each generator,
    the output displays the current representation, the generator index
    prefixed with \tt{g} for gauge or \tt{f} for fermion generators
    and followed by a $*$ suffix if the generator is $T$-odd,
    and then the actual generator definition.
\egroup

\item [\tt print geodesics $*$] \bgroup \itemsep 0pt \parskip 0pt
\item [\tt print geodesics {\it obs\#} {\tt{[}}{\it gen\#}{\tt{]}}]~\\
    Print symbolic geodesic equations.
    The first form prints all geodesic equations.
    The second form allows one to specify a particular observable
    and, optionally, the $T$-even vacuum representation generator
    which acts on the indicated observable.
\egroup

%\pagebreak[2]
\item [\tt print gradient $*$] \bgroup \itemsep 0pt \parskip 0pt
\item [\tt print gradient {\it term\#} {\tt{[}}{\it gen\#}{\tt{]}}]~\\
    Print symbolic gradient expressions.
    The first form prints all gradient components.
    The second form allows one to print the gradient components of a specified
    term  in the Hamiltonian or free energy
    (without any overall adjustable coupling constant) and,
    optionally, specify the index of the
    $T$-even vacuum representation generator which acts on the indicated term.
\egroup

\item [\tt print hamiltonian$^+$]~\\
    Print the expression defining the Hamiltonian.

\item [\tt print lagrange $*$] \bgroup \itemsep 0pt \parskip 0pt
\item [\tt print lagrange {\it $T$-even-gen\#} {\tt{[}}{\it $T$-odd-gen\#}{\tt{]}}]~\\
    Print symbolic Lagrange bracket expressions.
    The first form prints all Lagrange bracket components.
    The second form allows one to specify indices of one or both of the
    generators, in the active symmetry representation,
    whose commutator defines a Lagrange bracket component.
\egroup

\goodbreak

\item [\tt print mode$^+$ *]\bgroup \itemsep 0pt \parskip 0pt
\item [\tt print mode$^+$ {\tt{[}}{\it mode\#}{\tt{]}}]~\\
    Print all, or a specified, small oscillation eigenmode
    from the most recently evaluated spectrum.
\egroup

\item [\tt print observable $*$] \bgroup \itemsep 0pt \parskip 0pt
\item [\tt print observable {\it obs\#$_1$} \tt{[}{\it obs\#$_2$}\tt{]}]
\item [\tt print observable {\it obs$_1$} {\tt{[}{\it obs$_2$}\tt{]}} $\cdots$]
\item [\tt print observable ({\it obstype}) {\tt{[}\textit{...}\tt{]}}]
\item [\tt print observable ({\it corder},{\it xorder}) {\tt{[}\textit{...}\tt{]}}]
\item [\tt print observable ({\it corder},{\it xorder},{\it obstype}) {\tt{[}\textit{...}\tt{]}}]~\\
    Print observables.
    The first form prints the entire list of constructed observables.
    The second form allows one to 
    specify a single observable index, or a pair
    of indices indicating a range of observables.
    The third form allows one to input the printable form of one or more observables
    and determine if the given observables are present in the observable list.
    Canonicalization is not automatically performed on the input observables.
    The final forms allow one to print all observables of a specified type,
    with specified creation and expectation orders or, in the final form,
    specified orders and type.
    No spaces may be present within these parenthesized specifications.
    Multiple such specifiers may be given.
\egroup

\item [\tt print obsstats]~\\
    Print a table of observable statistics showing the breakdown of the
    number of constructed observables by type, creation order
    and expectation order.

\item [\tt print operator $*$] \bgroup \itemsep 0pt \parskip 0pt
\item [\tt print operator {\it op\#$_1$} {\tt{[}}{\it op\#$_2$}{\tt{]}}]~\\
    Print operators(s).
    The first form prints all defined operators (or ``\tt{Op}'s'')
    whose linear combinations, when symmetry projected, form generators.
    The second form allows one to 
    specify a single operator index, or a pair
    of indices to indicate a range of operators.
\egroup

\item [\tt print polystats]~\\
    Prints polynomial script statistics,
    including the number of polynomial terms and their breakdown
    by monomial order for the gradient, curvature and Lagrange bracket
    (in all representations), and geodesic equations.

\item [\tt print primary]~\\
    Print the list of primary \tt{Op}'s (those which cannot be
    produced by commutators of lower order \tt{Op}'s).

\item [\tt print representation] \bgroup \itemsep 0pt \parskip 0pt
\item [\tt print representation $*$]
\item [\tt print representation {\it irrep-name}]~\\
    Print irreducible representation (irrep) projectors.
    The first form lists all defined irrep names,
    while the second form prints the full form of all irreps.
    The final form allows one to specify a particular irrep by name,
    which is printed in full.
    For multi-dimensional representations, the \textit{irrep-name} actually
    specifies both the representation and a particular row of the
    representation matrices.
    
\egroup

\item [\tt print rkmethods]~\\
    Print the names of available Runge-Kutta ODE integration methods,
    and report the currently selected method.

\item [\tt print spectrum$^+$]~\\
    Prints the most recently evaluated small oscillation spectrum.

\item [\tt print stage]~\\
    Print the current minimization stage.

\item [\tt print state]~\\
    Print information on the theory class, lattice size,
    and global settings including
    minimization stage, active representation,
    ODE and minimization tolerances, ODE integration method,
    output file directories, etc.

\item [\tt print stats]~\\
    Print various operational statistics including the
    numbers of observables built, observable classifications,
    canonicalizations,
    approximations, symmetry transforms, etc., 
    numbers of minimizations and ODE integration steps,
    and canonicalization cache performance figures.

\item [\tt print symmetry $*$] \bgroup \itemsep 0pt \parskip 0pt
\item [\tt print symmetry {\it symmetry-name}]~\\
    Print symmetry transformations.
    The first form prints all symmetry transformations,
    while the second form allows one to specify a particular symmetry
    by name.
\egroup

\item [\tt print symmsets]~\\
    Print the various symmetry subsets used in observable canonicalization.

\item [\tt print sysindex]~\\
    Print a table showing the contents of the system information index.

\item [\tt print vevindex]~\\
    Print a list of data sets in a currently open vev-data file.

\item [\tt print version]~\\
    Print the program version.

\item [\tt purge observables]~\\
    Write, if needed, sys-info file(s) containing all observable information
    and then purge the master observable list, thereby freeing the memory
    which the list consumed.
    Any subsequent commands which require accessing the observable list will
    trigger reloading observable information into memory,
    but once all needed polynomial scripts have been generated
    numerical minimization and spectrum calculations may proceed
    without observable information in memory.

\item [\tt quit]~\\
    Exit the program.

\item [\tt read {\it filename}]~\\
    Read command input from the specified file.

\item [\tt reset blab]\bgroup \itemsep 0pt \parskip 0pt
\item [\tt reset maxthread]
\item [\tt reset mintol]
\item [\tt reset odetol]
\item [\tt reset representation]
\item [\tt reset rkmethod]
\item [\tt reset stage]
\item [\tt reset svdcutoff]
\item [\tt reset MMAobs]~\\
    Reset the indicated global variables to their default initial values.
    See the various \tt{set} commands for details.
\egroup

\item [\tt reset numerics]~\\
    Re-initialize observable expectation values
    for the current minimization stage.
    For gauge minimization, this zeros all pure gauge
    observable expectation values (except the identity).
    For fermi minimization, this zeros all fermion
    bilinear expectation values except those for which
    the fermion and conjugate fermion reside on the same
    site.  Such bilinears are initialized to their
    correct vacuum expectation values in the large mass limit.
    For Hamiltonian theories, this is minus half
    the value of the corresponding Wilson loop.
    For Euclidean theories, it is minus the inverse of the fermion mass
    times the value of the corresponding Wilson loop.

\item [\tt save sys]~\\
    Write system information for the current minimization stage
    (gauge or fermion) to a binary \tt{.sys} file.
    Sys-info files record the defined
    operators, observables, generators, and the polynomial scripts
    defining gradient, curvature, Lagrange bracket and geodesic equations,
    plus build-related statistics counters for the current minimization stage.
    Sys-info files are placed in theory-specific subdirectories
    (e.g., \tt{YM2e} or \tt{QCD2h})
    of the current save directory.

    Gauge information is written to files 
    named to indicate the Yang-Mills theory, including lattice dimension and whether
    the theory is Hamiltonian or Euclidean,
    the maximum order of gauge generators and observables, and
    a two-letter upper-case hash code of the gauge observable set.
    If observable approximation is turned on, an addition suffix
    staring with ``a'' is added to the file name.
    For example, \tt{YM2h_4g20QQ.sys} indicates a sys-info file for
    two-dimensional Hamiltonian Yang-Mills with order 4 gauge
    generators and order 20 gauge observables.

    Fermion information is written to files with names
    starting with \tt{QCD} followed by the theory specification and
    the maximum orders of generators and observables,
    and the associated observable set hash codes,
    for both the gauge and fermions stages.
    For example, \tt{QCD2h_4g20QQ_3f10CY.sys} denotes a fermion 
    sys-info file for two-dimensional Hamiltonian QCD with order 4 gauge
    generators and order 20 gauge observables together with
    order 3 generators and order 10 fermion observables,
    which is compatible with an associated \tt{YM2h_4g20QQ.sys}
    file of gauge information.

\item [\tt save vev]~\\
    Append to a binary \tt{.vev} data file the expectation
    values of all defined observables for the current minimization stage
    and current values of all adjustable coupling constants for the current stage.
    Multiple sets of such ``vev'' data may be present
    in the same \tt{.vev} file.
    The naming of vev-data files is identical to that of
    sys-info files, but with \tt{.vev} extensions.
    Vev-data files are placed, together with sys-info files,
    in theory-specific subdirectories of the current save directory.

\item [\tt set {\it coupling-name coupling-value}]~\\
    Sets the value of the named adjustable coupling constant.
    Changing the value of an adjustable coupling constant
    automatically changes the minimization stage to that
    associated with the coupling.

\item [\tt set approx true $\vert$ false]~\\
    Specifies whether non-retained observables are
    be approximated, when possible, by products of
    lower order retained observables (default \tt{false}).

\item [\tt set autoEgens true $\vert$ false]~\\
    Specifies whether 
    \tt{Eloop} generators are to be automatically
    created from each pure loop generator by
    evaluating the commutator of the loop with
    the gauge kinetic energy
    (default \tt{true}).

\item [\tt set autosave false $\vert$ true]~\\
    Specifies whether a sys-info file should automatically be
    written during geodesic script
    construction, and if a data record recording
    coupling values and observable vevs is automatically appended
    to a vev-data file after every successful minimization
    (default \tt{true}).

\item [\tt set blab {\it source-file verbosity-level}]~\\
    Sets the verbosity level for particular source files.
    The specified source file name should omit the
    extension, is case sensitive, and cannot be further abbreviated.
    Current settings are reported by \tt{print blablevels}.
    Other than \tt{set blab Build 1} or \tt{set blab Numerics 1},
    non-zero values are generally only relevant when
    debugging and can produce voluminous output.

\item [\tt set checkobs false $\vert$ true]~\\
    Specifies whether extra validity checking is performed
    on every constructed observable (default \tt{false}).

\item [\tt set dots false $\vert$ true]~\\
    Specifies whether observables are printed with
    periods separating the printed forms of individual symbols
    (default \tt{false}).

\item [\tt set echo false $\vert$ true]~\\
    Specifies whether command input lines are echoed
    to the output stream (default \tt{false}).

\item [\tt set gennorm false $\vert$ true]~\\
    Specifies whether to normalize generators
    as described in section 4.7 of the
    design and implementation [[notes]]
    (default \tt{false}).

\item [\tt set geoswap false $\vert$ true]~\\
    Specifies whether individual buckets of
    geodesic equation polynomials are to be
    swapped to and from disk as needed
    (default \tt{false}).
    Setting \tt{geoswap} true automatically
    also sets \tt{autosave} true.

\item [\tt set inputdir {\it directory}]~\\
    Specifies the direectory in which to find
    command input files without full pathnames.
    (default \tt{./}).

\item [\tt set maxnewt {\it integer}]~\\
    Specifies the maximum number of
    Newton iteration steps to allow
    during minimization
    (default 500).

\item [\tt set maxode {\it integer}]~\\
    Specifies the maximum number of
    ODE integration steps to allow
    during each geodesic equation integration
    (default 1500).

\item [\tt set maxthread {\it integer}]~\\
    Specifies the maximum number of
    threads to spawn during multi-threaded
    build or minimization stages.
    If zero (default), the number of threads is automatically
    controlled by \textsl{TBB} threading library.
    However, if \tt{geoswap} is set true, then
    the number of threads in use may be limited
    by a requirement that the per-thread I/O buffers not
    consume more than half of total system memory.

\item [\tt set mintol {\it value}]~\\
    Specifies the minimization convergence criterion
    (default $10^{-10}$ for 64-bit mode,
    $10^{-6}$ for 32-bit mode)).

\item [\tt set odetol {\it value}]~\\
    Specifies the ODE error tolerance
    (default $10^{-10}$ for 64-bit mode,
    $10^{-6}$ for 32-bit mode)).

\item [\tt set oknegeig false $\vert$ true]~\\
    Specifies whether the presence of negative
    curvature eigenvalues (for sectors other than 
    Euclidean fermions) will cause minimization to abort
    (default \tt{false}).

\item [\tt set representation {\it repname}]~\\
    Specifies the active symmetry representation.
    (\tt{print rep} lists the defined irreps.)

\item [\tt set rkmethod {\it method-name}]~\\
    Specifies the Runge-Kutta ODE integration method.

\item [\tt set savedir {\it directory}]~\\
    Specifies the directory in which to write
    system information and expectation value data
    files (default \tt{./save/}).
    Files are placed in (automatically created)
    theory-specific subdirectories of this directory.

\item [\tt set stage gauge $\vert$ fermi]~\\
    Specifies the minimization stage
    (initially set to \tt{gauge}).
    As noted above,
    changing the value of an adjustable coupling constant
    automatically changes the minimization stage to that
    associated with the coupling.

\item [\tt set svdcutoff {\it value}]~\\
    Sets the cutoff below which
    curvature singular values are truncated
    (default $0$).

\item [\tt set symcurv true $\vert$ false]~\\
    Specifies whether the curvature is
    symmetrized when evaluating the
    small oscillation spectrum
    (default \tt{true}).

\item [\tt set timing true $\vert$ false]~\\
    Specifies whether to report the time taken by each command,
    when greater than 10 milliseconds
    (default \tt{true}).

\item [\tt set vevappend false $\vert$ true]~\\
    When writing a vev-data file specifies if the output file,
    if pre-existing, is to be opened in append mode
    instead of being overwritten
    (default \tt{false}).  The setting applies to both
    gauge and fermion stage vev-data files.

\item [\tt set xtraobs false $\vert$ true]~\\
    Specifies whether observables which are generated
    during the building of curvature or Lagrange bracket scripts,
    but not already present in the set of built observables,
    are to be added to the list of observables.
    Such observables are not fully classified as to their
    creation and expectation orders.
    (default \tt{false}).

\item [\tt set MMAappend false $\vert$ true]~\\
    Specifies whether the output results file,
    if pre-existing, is to be opened in append mode
    instead of being overwritten
    (default \tt{false}).  The setting applies to both
    gauge and fermion stage MMA results files.

\item [\tt set MMAdir {\it directory}]~\\
    Specifies the directory in which to
    write output results in Mathematica entry format
    (default \tt{./MMA/)}.
    Files are placed in (automatically created)
    theory-specific subdirectories of this directory.

\item [\tt set MMAobs {\it obs$_1$} \tt{[}{\it obs$_2$}\tt{ [...]]}]~\\
    Appends the listed observables to the subset of observables whose
    expectation values are included in output written
    to the numerical \tt{.m} results file.
    The default set of such observables are those appearing
    in the Hamiltonian or free energy.

\item [\tt test arithmetic 32 \tt{|} 64]~\\
    Checks if the program is using 32 or 64 bit numbers
    to represent observable indices and polynomial coefficients.
    If the specified setting disagrees with the running program,
    an exception is raised and processing of any current command
    file is terminated.
    This command
    may be used within a command input file to signal that the file
    is designed for use with either 32 or 64 bit program versions.

\item [\tt test irreps]~\\
    Test the defined irreducible representation projectors for
    orthogonality, normalization, and completeness.
    (Used for debugging.)

\item [\tt test jacobi] \bgroup \itemsep 0pt \parskip 0pt
\item [\tt test jacobi {\it obs\#}]
\item [\tt test jacobi {\it op$_1$} {\it op$_2$} {\it obs\#}]~\\
    Test the validity of the Jacobi identity.
    The first form tests all triples of defined generators
    of the currently active representation.
    The second form tests the identity on the 
    observable with the specified index and all pairs
    of defined generators (of the currently active
    representation).
    The final form uses user-specified
    coherence algebra \tt{Op}'s and a specified observable.
    (Used for debugging.)
\egroup

\item [\tt write {\it filename}]~\\
    Redirect command output to the specified file.
    If the given filename is `\tt{-}', then subsequent
    output is returned to the initial output stream.

\end{description}

%\newpage

\begin{thebibliography}{99}

%\cite{YaffeRMP}
\bibitem{YaffeRMP}
L.~G.~Yaffe,
``Large $N$ limits as classical mechanics,''
\href{https://doi.org/10.1103/RevModPhys.54.407}
{Rev.~Mod.~Phys.~\textbf{54}, 407 (1982)}.
%doi:10.1103/RevModPhys.54.407
%305 citations counted in INSPIRE as of 01 Dec 2024

%\cite{CSVAI}
\bibitem{CSVAI}
F.~R.~Brown and L.~G.~Yaffe,
``The coherent state variational algorithm:
A numerical method for solving large-$N$ gauge theories,''
\href{https://doi.org/10.1016/0550-3213(86)90318-4}
{Nucl.~Phys.~B \textbf{271}, 267-332 (1986)}.
%doi:10.1016/S0550-3213(86)80012-8
%24 citations counted in INSPIRE as of 01 Dec 2024

%\cite{CSVAII}
\bibitem{CSVAII}
T.~A.~Dickens, U.~J.~Lindqwister, W.~R.~Somsky and L.~G.~Yaffe,
``The coherent state variational algorithm. 2. Implementation and testing,''
\href{https://doi.org/10.1016/0550-3213(88)90234-9}
{Nucl.~Phys.~B \textbf{309}, 1 (1988)}.
%doi:10.1016/0550-3213(88)90234-9
%20 citations counted in INSPIRE as of 01 Dec 2024

%\cite{Yaffe:2026ouz}
\bibitem{Yaffe:2026ouz}
L.~G.~Yaffe,
``Resurrecting the coherent state variational algorithm for large $N$ gauge theories,''
\href{https://arxiv.org/abs/2601.20150}
{arXiv:2601.20150 [hep-lat]}.
%0 citations counted in INSPIRE as of 30 Jan 2026

\end{thebibliography}
\end{document}
